% Transcription — fonds Alexandre Grothendieck, shelfmark 156-7, batch 3 (pages 41-60).
% Established from the facsimile archives/batches/156-7/batch-03.pdf (a copy of
% 156-7.pdf fetched through the project's relay; no cover sheet in the batch,
% so PDF page k = folder page 40+k), per the skill transcribe-grothendieck. The
% folder is 113 pages in six batches, transcribed in parallel; this is the
% third. The second draft (folder 156-6, GF VI) was read for notation, and
% the tail of batch 2 of this folder (pages 39-40, the restart with a dense
% part Λ: MΛ 0-MΛ 3, Lemmes A, B, C) once it existed; the other batches of
% this folder could not be relied on.
%
% Pass: Opus 5.5 (claude-opus-5-5), 2026-09-26 — first pass, unchecked against the pages by a human.
%
% Pages skipped: none. Pages summarised: none. All twenty leaves are his
% mathematics, in ink, fast drafting; all are transcribed.
%
% His pagination 41-60 stands at the head of each leaf and coincides with the
% archivists' (noted on page 41). Dates in his hand, each set as a
% \subsection with a \marginal and a \note where it stands: « 24 juin »,
% underlined, top left of page 45; « 25 juin », top left of page 54. Both
% fall inside the 23-26/06/1986 of the catalogue title.
%
% Numbered runs, recorded where they stand. Page 41: a struck NB, a margin
% « MΛ4 », Lemme A' (after Lemme A of p. 40), MΛ5 (the 5 over another
% figure), Lemme B (margin « Lemme 6 p. 32 »), a Cor, a cancelled Lemme C'.
% Page 42: a Proposition listing MΛ0-MΛ5 (margin names: axiome trivial, des
% raffinements induits, de densité de Λ, des strates minimales), and the
% relation 𝓡(X, Y). Page 43: At C_Λ (spac.) and a Cor. Page 44 (cancelled by
% three diagonal strokes, transcribed): MΛ'4, MΛ'5, margin « At D », a Lemme
% with conditions (i), [(ii)], (i'), (ii'). Page 45 (24 June): AtΛ1, AtΛ2,
% AtΛ3 « ou At DΛ » (margin: « si Λ = ℒ, c'est At L 3 »), At CΛ (spéc),
% At CΛ (pol), and a table of implications among Λ1-Λ3, C, C*, CΛ*, D, D_0
% (« cf. page 22 »). Page 46: four equivalent complete systems a)-d); the
% cases 1°) Λ = ℒ, 2°) Λ = 𝓜; « Λ-spécieux », « Λ-polyédral ». Page 47:
% MΛ0-MΛ3 « comme avant (p. 42) », MΛ4, MΛ5 (margin « remplace MΛ5 p. 34 »,
% margin « Non, cf. p. 49 »). Page 48: Lemme 1 a), b), a « Terminologie ».
% Page 49: Lemmes 2 and 3, and the admission « On tourne en rond ! ». Page 50:
% « Récapitulation » (his heading, set as \subsection; cites p. 23, p. 22,
% pp. 45-46). Page 51: three conceptions (1)-(3) (circled numerals, set in
% parentheses): Prat 1-Prat 6 (no Prat 5; « Prat 4 » over a « Prat 5 »),
% PratΛ1-PratΛ6. Page 52: (3) Prat dis 1'-Prat dis 6' (« dis » struck and
% replaced by ℒ in 2'-6'), cf. p. 34. Page 53: Prat pol, a Remarque (the
% « trivial » préatelier, ≪ = ≤), Prat 1 - Prat 6, At D, At C (spéc), At Fil,
% At supp, At L 1-3. Page 54 (25 June): Proposition with conditions (i)
% (a), b), c)) and (ii), citing « GF VI p. 51 » and « p. 40 », « p. 45 ».
% Pages 55-58: its proof, Lemme 1 (φ_Λ), an Exemple (affine simplices in a
% space X) and a Cor; a counter-example (Figpol(ℒ), p. 58), Prat 1' - Prat 6'
% cited. Page 59: Déf of « ensembliste spécieux / polyédral » and
% « quasi-ensembliste » ateliers; Prat quens 1-4. No formula numbers (1.x) in
% this batch. Internal cross-references, kept as plain page numbers: p. 30,
% p. 32 (p. 41), p. 22, p. 23 (pp. 45, 50), p. 34 (pp. 47, 52), p. 42
% (pp. 47, 49), p. 49 (p. 47), p. 6 (p. 51), pp. 40, 45 (p. 54).
% Cross-references to the second draft: « GF VI p. 51 » (p. 54, plongements
% d'ateliers), noted there; notes on p. 45 (AtΛ1 generalises At 8, GF VI
% p. 45, (1.52)) and p. 54 (the ensemblist ateliers of GF VI pp. 52-58).
%
% What the batch is about. Continuing from the dense part Λ ⊂ 𝓜 introduced
% at the end of batch 2, he tries to characterise the systems
% (𝓜, ≤, ≪, compatibility, Λ) that come from an atelier (Proposition p. 42,
% axioms MΛ0-MΛ5: trivial axiom, induced refinements, density of Λ, minimal
% strata, saturation of compatibility, and MΛ5 in terms of the relation 𝓡),
% then to define compatibility from a disjointness relation on Λ alone
% (At C_Λ (spac.), MΛ'4-5, pp. 43-44, cancelled). On 24 June (p. 45) he
% sets out AtΛ1-3 and At CΛ (spéc/pol), charts their implications, and names
% « Λ-spécieux » and « Λ-polyédral » ateliers (pp. 45-46). Pages 47-49 try
% again with MΛ4-5 and Lemmes 1-3 and end: « Rien de nouveau ! On tourne en
% rond ! ». Page 50 recapitulates and proposes (𝓜, ≤, ≪, compatibility) as
% an independent structure, a « pré-atelier ». Pages 51-53 give three
% axiomatics: (1) Prat 1-6, (2) PratΛ1-6, (3) Prat dis 1'-6' in terms of a
% disjunction on ℒ, plus Prat pol for the polyhedral case, and the
% « trivial » préatelier (≪ = ≤), which satisfies everything except
% divisibility. On 25 June (pp. 54-58) he returns to ensemblist ateliers:
% the canonical φ: 𝔉 → Fig(ℒ), F ↦ multomb(F); (i) φ an embedding onto a
% spécieux sub-atelier vs (ii) 𝔉 spécieux with distinct lieux non-disjoint;
% (i) ⇒ (ii), (ii) ⇒ (i a) via Lemme 1 on φ_Λ; injectivity fails (affine
% structures on a segment), and a polyhedral counter-example for b). Page 59
% defines « ensembliste spécieux/polyédral » and « quasi-ensembliste »
% ateliers and gives Prat quens 1-4; page 60 reflects on why this common
% structure suits both contexts and on « admissible » closed parts defined by
% local finiteness conditions.
%
% Notation, aligned with GF VI and batch 2: 𝔉 \mathfrak{F}, 𝓜 \mathcal{M}, ℒ
% \mathcal{L}, Λ \Lambda, compatibility \between, disjointness \parallel,
% ≪ \ll, interior refinement \overset{\circ}{\ll}, « omb », « Omb »,
% « multomb », « Fig », « Figpol », « Ssfig », « Inf » in roman; « MΛ »,
% « Prat », « At » labels as he writes them; a
% vertical ≪ with X' written under X is rendered as X' ≪ X with a note; his
% underlined ≪ (pp. 41, 49, 51) is \underline{\ll}; 𝓡 \mathcal{R} for his
% round R (p. 42); 𝒜 \mathcal{A} for the atelier (pp. 46, 54); Φ for a
% set-theoretic figure (p. 56); the negated ∥ (p. 54) \overline{\parallel}.
% His underlining is \emph; circled (1)-(3) are set in parentheses.
%
% Hard pages: 53 (bottom-left margin a tangle of strikes), 56-57 (the
% injectivity argument and its example, much abbreviated), 58 (a cancelled
% paragraph), and above all 60, a page of fast prose of which many words are
% left \ill{}. Readings to check first: the margin notes of pp. 42, 51, 53,
% 59, 60; « prédisant » (p. 53); « faux » (p. 58); the phrase « (i a, c)
% non : ni X ≬ Y » (p. 57); the first line of p. 56. Variances on the page
% kept and flagged: {X, Y} then Inf(Y, Z) in Prat pol (p. 53); PratΛ6 with
% « Z' ≪ Z ≤ X, Y » against MΛ5's Z ≪ Z' (p. 51); « omb(Y') » for
% omb(Y_{X'}) (p. 59); « mais X ≬ Y » where non-compatibility is meant
% (p. 56).
\documentclass[11pt,a4paper]{article}
\input{../preamble/grothendieck}

\folder{156-7}
\batch{3}
\pages{41}{60}
\dating{1986}
\watermark{Édition de démonstration}
\foldertitle{[Chapitre] VII. Analysis situs (troisième mouture) : notes manuscrites (23-26/06/1986)}

\begin{document}
\page{41}
\note{pagination de sa main 41 à 60 en tête des feuillets, qui coïncide avec celle des archivistes}

\struck{NB Cet axiome ne fait intervenir $\between$ que par l'intermédiaire de la relation $\parallel$ ; et même seulement de la relation $X' \parallel Y'$ entre éléments de $\Lambda$.}
\note{ce NB, en tête de page, est barré de deux longs traits obliques}

\marginal{\emph{M$\Lambda$4} $X \between Y$, $X' \leq X$, $Y' \leq Y$ \uncertain{$\Longrightarrow$} $X' \between Y'$}
\note{note oblique dans la marge gauche, en face du NB biffé ; « M$\Lambda$4 » est souligné}

\textbf{\emph{Lemme} A'} Comme au lemme A, on suppose seulement $X \between Y$.

Mais on a besoin dans le raisonnement (p.~30) de la relation
\[
  \mathrm{omb}_{\Lambda}(X) \cap \mathrm{omb}_{\Lambda}(Y) = \mathrm{omb}_{\Lambda}(L) \overset{\mathrm{def}}{=} \bigcup_{Z \in \widetilde{X} \cap \widetilde{Y}} \mathrm{omb}_{\Lambda} Z
\]
ce qu'il faut \uncertain{poser} ici en axiome :

\textbf{\emph{M$\Lambda$5}} Si $X \between Y$, et si $Z \ll X$, $Z \ll Y$, alors $\exists\, Z' \in \widetilde{X} \cap \widetilde{Y}$ avec $Z \underline{\ll} Z'$, i.e.\ $Z \ll Z' \leq X, Y$.
\note{le « 5 » de « M$\Lambda$5 » est écrit en surcharge sur un autre chiffre (un 4 ?) ; l'énoncé est marqué d'un trait vertical dans la marge ; le signe entre $Z$ et $Z'$ est un $\ll$ souligné, que l'« i.e. » explicite}
\marginal{On peut se borner au cas $Z \in \Lambda$}

(Propriété \uncertain{restrictive} sur les \struck{\uncertain{couples}} \add{paires} $(X, Y)$ avec $X \between Y$)

\textbf{\emph{Lemme} B} Soient $Z \overset{\circ}{\ll} X$, $Z \overset{\circ}{\ll} Y$, avec $X \between Y$. Alors $X = Y$ (i.e.\ \uncertain{c'est} $\mathrm{At}_{\mathcal{M}}$ 4 b))
\note{le « B » est écrit en surcharge ; dans la marge gauche, « \uncertain{Lemme} 6 p.~32 »}

\textbf{\emph{Cor}} Plus gén\supplied{éralement}, si \struck{$Z \overset{\circ}{\ll} X$, $Z \ll Y$,} $X \between Y$ \add{et} $Z \overset{\circ}{\ll} X$, \struck{donc} $Z \ll Y$, alors $X \leq Y$. \struck{Utilisant M$\Lambda$2 (i.e.\ $\exists\, Z_0 \in \Lambda$, $Z_0 \ll Z$ donc $Z_0 \ll X$) et $Z_0 \ll Y$, \ill{} \ill{} condition \ill{}}

\struck{\emph{Lemme} C' Soient $X$, $Y$ avec \ill{} $X \between Y'$. Conditions équivalentes : (i) $Y \ll X'$ \ill{}} \marginal{\struck{Conditions équivalentes et $Z \ll X$.}}
\note{le lemme C' est annulé par trois traits obliques et plusieurs biffures ; lecture très partielle ; un « (\ill{}) » sous le bloc biffé, à gauche de « alors », n'a pas été lu}

En effet, soit $Z_0 \in \Lambda$, $Z_0 \overset{\circ}{\ll} Z$, alors $Z_0 \overset{\circ}{\ll} Z \overset{\circ}{\ll} X$ donc $Z_0 \overset{\circ}{\ll} X$, \struck{\ill{}} i.e.\ $Z_0 \in \mathrm{omb}_{\Lambda}(X)^{\circ}$, $Z_0 \ll Z \ll Y$ donc $Z_0 \ll Y$, i.e.\ $Z_0 \in \mathrm{omb}_{\Lambda}(Y)$, donc $\mathrm{omb}_{\Lambda}(X)^{\circ} \cap \mathrm{omb}_{\Lambda}(Y) \neq \emptyset$, donc $X \leq Y$ par lemme A'. OK.

\page{42}
\textbf{\emph{Proposition}} Propriétés caractéristiques d'un système $(\mathcal{M}, \leq, \ll, \between, \Lambda \subset \mathcal{M})$, pour provenir d'un atelier avec $\Lambda$ dense dans $\mathcal{M}$.

\textbf{M$\Lambda$0} $X \leq Y \Longrightarrow X \ll Y$
\marginal{Axiome trivial}

\textbf{M$\Lambda$1} Si $Y \ll X$, $X' \ll X$, $\underline{Z \in \Lambda}$ et $Z \ll Y$, $Z \ll X'$, alors $\exists\, Y'$ avec $Y' \leq Y$, $Y' \ll X'$, $Z \ll Y'$ :
\[
  \mathrm{omb}_{\Lambda}(Y) \cap \mathrm{omb}_{\Lambda}(X') = \bigcup_{Y' \in \struck{\widetilde{Y} \cap \mathrm{Omb}(X')}} \mathrm{omb}_{\Lambda}(Y')
  \qquad \overset{\text{si } Y \leq X}{=} \qquad \bigcup_{Y' \in \widetilde{Y} \cap \widetilde{X}'} \mathrm{omb}_{\Lambda}(Y')
\]
\note{la page écrit $X'$ sous $X$, relié par un $\ll$ vertical dont la pointe va vers $X'$, lu $X' \ll X$ ; le « $X$ » de $X'$ est repassé. « $Z \in \Lambda$ » est souligné d'un trait ondulé. Sous l'indice biffé de la première réunion, une flèche désigne l'« $\mathrm{Omb}$ » : « grande ombre »}
\marginal{Axiome des raffinements induits}
\marginal{\struck{Mais seul} \ill{} \ill{} M$\Lambda$4 + M$\Lambda$5 \ill{}}
\note{deux notes obliques dans la marge gauche, à la hauteur de M$\Lambda$0 et de M$\Lambda$1 ; la seconde est en partie biffée et n'est lue que par fragments}

\textbf{M$\Lambda$2} $\forall\, X \in \mathcal{M}$, $\exists\, Z \in \Lambda$ avec $Z \overset{\circ}{\ll} X$
\marginal{Axiome de densité de $\Lambda$}

\textbf{M$\Lambda$3} $\forall\, X \in \mathcal{M}$, $Z \in \Lambda$, avec $Z \ll X$, $\exists\, X' \leq X$ avec $Z \overset{\circ}{\ll} X'$
\marginal{Axiome des strates minimales}

\textbf{M$\Lambda$4} \add{$\between$ relation réflexive symétrique, et} ($X \between Y$, $X' \leq X$, $Y' \leq Y \Longrightarrow X' \between Y'$)
\note{l'ajout est entouré et rattaché à l'énoncé par un trait}

\textbf{M$\Lambda$5} Si $X \between Y$ et $Z \in \Lambda$ tel que $Z \ll X$, $Z \ll Y$, $\exists\, Z' \in \mathcal{M}$ avec $Z \ll Z' \leq X, Y$, i.e.
\[
  \mathrm{omb}_{\Lambda}(X) \cap \mathrm{omb}_{\Lambda}(Y) = \bigcup_{Z' \in \widetilde{X} \cap \widetilde{Y}} \mathrm{omb}_{\Lambda}(Z')
\]
\note{le « 5 » est écrit en surcharge (sur un 4 ?)}
\marginal{Cette condition \ill{} signifie que \ill{} \ill{} \ill{} W \ill{} i.e.\ \ill{} fait \ill{} \ill{} pour $Z \in \Lambda$.}
\note{note oblique dans la marge gauche, en face de M$\Lambda$5, lue par fragments}

C'est \struck{des} \add{deux} derniers axiomes seulement qui fassent intervenir $X \between Y$. L'axiome M$\Lambda$4 est un axiome de saturation. Pour M$\Lambda$5, \struck{soit} Notons $\mathcal{R}$ la relation (définie en termes de $(\mathcal{M}, \leq, \ll)$ seulement, sans \uncertain{y} \ill{} $\Lambda$ ni $\between$) :
\[
  \mathcal{R}(X, Y) \iff
  \begin{array}{l}
    \forall\, Z \in \mathcal{M} \text{ avec } Z \ll X,\ Z \ll Y, \\
    \exists\, Z' \in \mathcal{M} \text{ avec } Z \ll Z' \leq X, Y \\
    \text{i.e. } \mathrm{Omb}(X) \cap \mathrm{Omb}(Y) = \bigcup_{Z' \in \widetilde{X} \cap \widetilde{Y}} \mathrm{Omb}(Z')
  \end{array}
\]
\note{le signe de la relation est une ronde « $\mathcal{R}$ » }

\page{43}
\add{Cette relation est réflexive symétrique.}
\note{ajout en tête de page, au-dessus de la première ligne}

On vérifie tout de suite (en utilisant M$\Lambda$1 dans le cas de $Y \leq X$, $X' \ll X$) que \struck{cette relation} si $X' \leq X$, $Y' \leq Y$, alors $\mathcal{R}(X, Y) \Longrightarrow \mathcal{R}(X', Y')$. Et l'axiome M$\Lambda$5 est \ill{} \ill{} pour cette condition. Ainsi l'axiome M$\Lambda$5 prend la forme
\[
  X \between Y \Longrightarrow \mathcal{R}(X, Y)
\]
et on voit donc que les relations \add{réfl.\ sym.} $\between$ sont les relations (« saturées pour $\leq$ ») qui impliquent $\mathcal{R}(X, Y)$. La relation admissible $\between$ la plus large qu'on \uncertain{obtienne}, on prendrait $\mathcal{R}$ elle-même.
\note{dans la première phrase, $X'$ est écrit sous $X$, relié par un $\ll$ vertical, comme page 42 ; la seconde mention de l'axiome s'écrit « ML 5 » sur la page, la troisième « M$\Lambda$5 »}

Dans le cas \add{d'un atelier} ensembliste spécieux, on a $\between = \mathcal{R}$, mais pas dans d'autres cas intéressants. Ainsi on note que $\mathcal{R}(x, y)$ toujours vrai si $x, y$ sont dans $\mathcal{L}$ — or dans les cas non ensemblistes, il \uncertain{pourrait} exister des lieux $x, y$ qui ne soient compatibles.

Je voudrais maintenant voir que la condition $\between$ peut se définir par la connaissance de $\parallel$ sur $\Lambda$, par une condition

\textbf{\emph{At C$_{\Lambda}$} (spac.)} \struck{Soient} $X, Y \in \mathcal{M}$, \struck{Alors $X \between Y$} $L = \widetilde{X} \cap \widetilde{Y}$.
\[
  X \between Y \iff
  \begin{array}{l}
    \forall\, X' \in \mathrm{omb}_{\Lambda} X,\ Y' \in \mathrm{omb}_{\Lambda} Y \text{ avec } X'_L = Y'_L = \emptyset , \\
    \text{on a } X' \parallel Y'
  \end{array}
\]
où $X'_L = \widetilde{X}' \cap \mathrm{Omb}(L) = \lbrace Z \mid Z \leq X',\ \exists\, S \in L,\ Z \ll S \rbrace$.
\note{la parenthèse de l'étiquette « (spac.) » est écrite sur un mot biffé ; sous l'équivalence, trois lignes biffées : « $\forall\, X', Y' \in \Lambda$, $X' \ll X$, $Y' \ll Y$ », « $\forall\, X' \in \mathrm{omb}_{\Lambda}(X) \cap \mathrm{omb}_{\Lambda}(L)$ », « $Y' \in \mathrm{omb}_{\Lambda}(Y) \smallsetminus \mathrm{omb}_{\Lambda}(L)$ », et une formule biffée « $X'_L = \mathrm{omb}_{\Lambda}(L) \cap \bigcup_{Z \in L} \ldots$ » ; l'énoncé est marqué d'un trait vertical dans la marge}

(Ceci est vrai, \struck{\ill{}}

\textbf{\emph{Cor}} Soient $X, Y \in \mathcal{M}$. \struck{Alors $X'$ tel que} \add{Alors} $X \parallel Y$ (i.e.\ \struck{$X \between Y$ et $L = \emptyset$}) \struck{si $X'$} \supplied{si} $\forall\, X' \in \mathrm{omb}_{\Lambda}(X)$, $\forall\, Y' \in \mathrm{omb}_{\Lambda} Y$ on a $X' \parallel Y'$.

\page{44}
\note{toute la page est barrée de trois longs traits obliques ; elle est transcrite, la page 47 reprend M$\Lambda$4 et M$\Lambda$5}

La question, c'est quand un système $(\mathcal{M}, \leq, \ll, \Lambda \subset \mathcal{M}, \parallel_{\Lambda})$ \add{avec $\Lambda \subset \mathcal{M}$ dense,} provient d'une structure d'\struck{A}telier satisfaisant \struck{At C$_{\Lambda}$ (spac.)} \add{At D}. (\struck{Cette} Les axiomes M$\Lambda$0 à M$\Lambda$3 restent \uncertain{en vigueur}.) \struck{Pour} On doit avoir
\marginal{\emph{At D} \ill{} \ill{}}
\note{la note marginale est encadrée d'un trait courbe, dans l'angle supérieur gauche ; « At D » est écrit au-dessus de « At C$_{\Lambda}$ (spac.) », qui n'est pas biffé nettement ; après « On doit avoir », une coche}

\textbf{\emph{M$\Lambda$'4}} \add{La relation} $\parallel$ \add{sur $\Lambda$} est symétrique et antiréflexive ($X \parallel Y$ implique $X \neq Y$). \struck{$Z \parallel Z'$} Si $X \parallel Y$ et $X' \ll X$, $Y' \ll Y$, alors $X' \parallel Y'$.

\textbf{\emph{M$\Lambda$'5}} Soient $Y \leq X$, \add{$X, Y \in \mathcal{M}$ tels qu'il existe $Z$ avec $X, Y \leq Z$.} \struck{\ill{} $L = \widetilde{X} \cap \widetilde{Y}$, \ill{} $X', Y' \in \Lambda$ \ill{} $X' \ll X$, $Y' \ll Y$}
$Y' \in \mathrm{omb}_{\Lambda}(Y)$, $X' \in \mathrm{omb}_{\Lambda}(X)$. Supposons $X'_Y = \emptyset$ [i.e.\ $\mathrm{omb}_{\Lambda}(Y) \cap \mathrm{omb}_{\Lambda}(X') = \emptyset$, ou encore \struck{$\widetilde{X}' \cap$} $\mathrm{Omb}_{\Lambda}(Y) \cap \widetilde{X}' = \emptyset$ ou encore $\mathrm{Omb}(Y) \cap \mathrm{Omb}(X') = \emptyset$ — tout se ramène au $=$]. Alors $X' \parallel Y'$.
\note{sous « M$\Lambda$'5 Soient », deux lignes sont couvertes de biffures et de gribouillis ; l'ajout « $X, Y \in \mathcal{M}$ tels qu'il existe $Z$ avec $X, Y \leq Z$ » est écrit au-dessus d'elles ; l'énoncé est marqué d'un trait vertical dans la marge}
\marginal{NB On définit $X'_Y = \widetilde{X}' \cap \mathrm{Omb}(Y)$}

Vérifions d'abord les équivalences entre \ill{} :

\textbf{\emph{Lemme}} (\struck{\ill{}} indépendant de la donnée de $\parallel$ et de $\between$ — ne dépend que de $\leq$, $\ll$, $\Lambda$ satisfaisant M$\Lambda$0 à M$\Lambda$3) Soient $Y \leq X$, $X' \ll X$. \struck{Conditions} Conditions équivalentes \add{(i), (i'), (ii') et impliquant (ii)} :
\begin{itemize}
  \item[(i)] $\mathrm{omb}_{\Lambda}(Y) \cap \mathrm{omb}_{\Lambda}(X') = \emptyset$
  \item[(ii)] $\mathrm{omb}_{\Lambda}(Y) \cap \widetilde{X}' = \emptyset$, $\longleftarrow$ est impliqué par les autres, et leur est équivalent si on suppose $\forall\, X \in \mathcal{M}$, $\exists\, Z \in \Lambda$ tel que $Z \leq X$ ($\Lambda$ « \emph{très} dense » dans $\mathcal{M}$)
  \item[(i')] $\mathrm{Omb}_{\Lambda}(Y) \cap \mathrm{Omb}(X') = \emptyset$
  \item[(ii')] $\underbrace{\mathrm{Omb}(Y) \cap \widetilde{X}'}_{X'_Y} = \emptyset$
\end{itemize}
\note{« $X' \ll X$ » : $X'$ est écrit sous $X$, avec un $\ll$ vertical ; la condition (ii) est entre crochets sur la page}

\struck{Dém (i) $\Longrightarrow$ (ii) tautologique}, (ii') $\Longrightarrow$ (i) : \struck{En} Supposons (ii), prenons (i), par M$\Lambda$1 on a $\mathrm{omb}_{\Lambda}(Y) \cap \mathrm{omb}_{\Lambda}(X') = \bigcup_{Y' \in \widetilde{X}' \cap \ldots}$
\note{la page s'arrête sur cette réunion, dont l'indice est coupé au bas du feuillet}

\page{45}
\subsection{24 juin}
\marginal{24 juin}
\note{date de sa main, soulignée d'un trait oblique, dans l'angle supérieur gauche de la page 45, en tête du premier alinéa}

Soit $\Lambda \subset \mathcal{M}$, fermé pour $\leq$. On introduit dans \uncertain{maintenant} des conditions analogues pour $\mathcal{L}$ :

\textbf{\emph{At$\Lambda$1}} $\forall\, X \in \mathcal{M}$, $\mathrm{omb}_{\Lambda}(X)^{\circ} \neq \emptyset$, i.e.\ $\exists\, Z \in \Lambda$, $Z \overset{\circ}{\ll} X$
\marginal{densité de $\Lambda$}
\note{pour $\Lambda = \mathcal{L}$, At$\Lambda$1 est l'« axiome des lieux » At 8 de la deuxième mouture (GF VI, p.~45, formule (1.52) $\mathrm{omb}(X)^{\circ} \neq \emptyset$), repris ici pour une partie $\Lambda$ quelconque}

\textbf{\emph{At$\Lambda$2}} $\forall\, X, Y \in \mathcal{M}$, $X \parallel Y \iff \forall\, X' \in \mathrm{omb}_{\Lambda}(X)$, $Y' \in \mathrm{omb}_{\Lambda}(Y)$, $X' \parallel Y'$
\marginal{si on a At$\Lambda$1 \uncertain{et} At$\Lambda$2, on dit que $\Lambda$ \uncertain{est} \uncertain{très} dense}
\note{« At$\Lambda$2 » est doublement souligné ; la note marginale, oblique, n'est lue qu'en partie}

\textbf{\emph{At$\Lambda$3}} ou \emph{At D$\Lambda$} Soient $X, Y \in \mathcal{M}$ \struck{tels qu'il existe} \struck{$X \between Y$,} \struck{\ill{}} $L = X \cap Y$, $X' \in \mathrm{omb}_{\Lambda}(X)$, $Y' \in \mathrm{omb}_{\Lambda}(Y)$. Alors
\[
  X'_L \parallel Y'_L \Longrightarrow X' \parallel Y' .
\]
\marginal{si $\Lambda = \mathcal{L}$, c'est At L 3}
\note{« At$\Lambda$3 » est doublement souligné ; sous $\parallel$ un signe noirci ; après « $L = X \cap Y$ », « déf » surmonte un mot biffé ; au-dessus, « Supposons que », biffé}

\textbf{\emph{At C$\Lambda$ (spéc)}} Soient $X, Y \in \mathcal{M}$, $L = X \cap Y$ (défini par $L = \widetilde{X} \cap \widetilde{Y}$). Alors
\[
  X \between Y \iff \forall\, X' \in \mathrm{omb}_{\Lambda}(X),\ Y' \in \mathrm{omb}_{\Lambda}(Y),\ X'_L = Y'_L = \emptyset \Longrightarrow X' \parallel Y'
\]

\textbf{\emph{At C$\Lambda$ (pol)}} Soient $X, Y \in \mathcal{M}$, $L = X \cap Y$. Alors
\[
  X \between Y \iff
  \begin{cases}
    \text{(i) Comme dans At C$\Lambda$ (spéc) ci-dessus} \\
    \text{(ii) $L = \emptyset$ ou $L \in \mathcal{M}$}
  \end{cases}
\]

On a en tout cas les implications (cf.\ page 22)
\[
  \begin{array}{ccccc}
    \Lambda 3 + C\Lambda^{*} & \Longrightarrow & C\,\struck{\ill{}} & & \\
     & & \Downarrow & & \\
    \Lambda 2 + \Lambda 3 & \Longrightarrow & D & \Longrightarrow & \Lambda 3 \\
    \Downarrow & & \Downarrow & & \\
    \Lambda 2 & \Longrightarrow & D_0 & &
  \end{array}
\]
\[
  D_0 + \Lambda 1 + C\Lambda^{*} \Longrightarrow \Lambda 2 , \qquad
  \Lambda 2 + C\Lambda^{*} \iff \Lambda 2 + C^{*}
\]
où on a laissé tomber les sigles « At » dans la désignation des axiomes, et où $C\Lambda^{*}$ désigne indifféremment At C$\Lambda$ (spéc) ou At C$\Lambda$ (pol), et \struck{C \ill{} désigne At C (\ill{} \ill{} \ill{})} $C^{*}$ est At C (spéc) ou At C (pol) (mais toujours les \uncertain{mêmes}).
\note{dans le tableau, la flèche verticale sous $\Lambda 2 + \Lambda 3$ est barrée d'un trait qui la croise (une implication annulée ?) ; la troisième ligne, $D_0 + \Lambda 1 + C\Lambda^{*} \Longrightarrow \Lambda 2$, aboutit sur la page au même $\Lambda 2$ que la colonne de gauche ; à la place de la lettre après « C » dans la première ligne, un griffonnage. Le renvoi « p. 22 » est à sa pagination ; l'ajout interlinéaire au-dessus de la parenthèse biffée n'est pas lu}

\page{46}
Un ensemble \emph{complet} de conditions de compatibilité disjonction impliquant $\Lambda$, est (impliquant donc tous les axiomes pertinents
\[
  \underset{\substack{\Downarrow \\ D_0}}{C,\ D},\ C^{*},\ \Lambda 1,\ \Lambda 2,\ \Lambda 3,\ C\Lambda^{*}
\]
pour \uncertain{la} conclusion) des suivants, p.\ ex.
\[
  \begin{array}{ll}
    a) & \Lambda 1,\ \Lambda 2,\ \Lambda 3,\ C\Lambda^{*} \\
    b) & \Lambda 1,\ \Lambda 2,\ D,\ C^{*} \\
    c) & \Lambda 1,\ C\Lambda^{*},\ D \\
    d) & \Lambda 1,\ \Lambda 2,\ C^{*},\ D_0
  \end{array}
\]
\note{la double flèche verticale « $\Downarrow D_0$ » est écrite sous le « $D$ » de la première liste ; les quatre listes a)-d) sont réunies par une accolade à gauche ; entre a) et b), une double flèche verticale relie $\Lambda 3$ à $D$, une autre, à double pointe, $C\Lambda^{*}$ à $C^{*}$ ; en d), « $\Lambda 2$ » et « $C^{*}$ » sont repassés}

les équivalences marquées valent toutes, en présence de $\Lambda 1$, $\Lambda 2$ i.e.\ si $\Lambda$ est partie très dense de $\mathcal{M}$.

\marginal{NB $\Lambda 1$ est le seul des axiomes pertinents qui figure dans chacun des ensembles (a), b), c), d)) \uncertain{équivalents}}
\note{note à droite de la liste d), en trois lignes serrées, lue en partie}

On \uncertain{dira}, quand cet ensemble de conditions est vérifié, qu'on a un atelier $\Lambda$-spécieux, ou $\Lambda$-polyédral. Les cas les plus intéressants pour moi sont
\begin{itemize}
  \item[1°)] $\Lambda = \mathcal{L}$ cas déjà traité — mais cela signifie que l'atelier satisfait \uncertain{maintenant} les conditions de divisibilité.
  \item[2°)] $\Lambda = \mathcal{M}$, auquel cas $\Lambda 1$ devient tautologique, $\Lambda 2$ équivaut à $D_0$, $\Lambda 3$ équivaut à $D$, $C\Lambda^{*}$ équivaut à $C^{*}$, \struck{et \ill{}} \ldots\ \emph{$\mathcal{M}$-polyédral} signifie. Dire que $\mathcal{A}$ est $\mathcal{M}$-spécieux signifie simplement qu'il satisfait : $D$, et $C$ spéc resp.\ $C$ pol.
\end{itemize}
\note{« $\mathcal{M}$-polyédral » est souligné ; un trait le relie à « signifie » à la ligne suivante}

Dans le cas général, on peut dire que $\mathcal{A}$ est (\uncertain{lorsqu'il est}) $\Lambda$-spécieux, dès \uncertain{qu'il} est $\mathcal{M}$-spécieux i.e.\ spécieux tout court, et si de plus $\Lambda$ est très dense ($\Lambda 1$ et $\Lambda 2$).
\note{$\mathcal{A}$ désigne ici l'atelier ; la lettre est une ronde « A »}

\page{47}
\note{le « 7 » de son numéro est repassé}
On voudrait donner une condition sur
\[
  (\mathcal{M},\ \underset{\mathcal{M}}{\leq},\ \underset{\mathcal{M}}{\ll},\ \Lambda \subset \mathcal{M},\ \underset{\Lambda}{\parallel})
\]
où $\underset{\Lambda}{\parallel}$ est une relation dans $\Lambda$, \add{partie fermée de $\mathcal{M}$,} pour que ce proviennent d'un atelier $\Lambda$-spécieux, resp.\ $\Lambda$-pol. On se \uncertain{propose} pour fixer les idées, prenons le cas $\Lambda$-spécieux.

\emph{M$\Lambda$0}, \emph{M$\Lambda$1}, \emph{M$\Lambda$2}, \emph{M$\Lambda$3} comme avant (p.~42)

\textbf{\emph{M$\Lambda$4}} : La relation $\underset{\Lambda}{\parallel}$ (ou $\parallel$) sur $\Lambda$ est \struck{réflexive} symétrique, antiréflexive ($X \parallel Y \Rightarrow X \neq Y$), et $X \parallel Y$, $X' \ll X$, $Y' \ll Y \Longrightarrow X' \parallel Y'$.

\textbf{\emph{M$\Lambda$5}} Soient $X, Y \in \mathcal{M}$ tels qu'il existe $Z \in \mathcal{M}$ avec $X, Y \leq Z$, \struck{$L = X \cap$} $X' \in \mathrm{omb}_{\Lambda}(X)$, $Y' \in \mathrm{omb}_{\Lambda}(Y)$. Posons
\[
  \begin{cases}
    L = \widetilde{X} \cap \widetilde{Y} \\
    X'_L = \widetilde{X}' \cap \mathrm{Omb}(L) , \quad Y'_L = \widetilde{Y}' \cap \mathrm{Omb}(L) \\
    \text{où } \mathrm{Omb}(L) \overset{\mathrm{def}}{=} \bigcup_{Z \in L} \mathrm{Omb}(Z)
  \end{cases}
\]
Supposons $X'_L \parallel Y'_L$ au sens : $\forall\, X'' \in X'_L$, $Y'' \in Y'_L$, \add{on a} $X'' \parallel Y''$ (NB on aura $X'' \leq X'$, $Y'' \leq Y'$, donc $X'', Y'' \in \Lambda$). Alors $X' \parallel Y'$.
\note{sous « $\mathrm{Omb}$ » de $X'_L$, un indice noirci ; l'énoncé est marqué d'un trait vertical dans la marge}
\marginal{remplace M$\Lambda$5 p.~34}
\marginal{On a besoin \ill{} \ill{} plus \ill{} \ill{} $X \neq Y$ (car \ill{} \ill{} dense).}
\note{deux notes obliques dans la marge gauche, en face de M$\Lambda$5 ; la seconde n'est lue que par fragments}

\uncertain{Plus} \add{bas} je dis que ça suffit (\add{plus} \uncertain{évidemment} \ill{} nécessaire). \struck{\ill{} \ill{}} On pose, pour $X, Y \in \mathcal{M}$, si $L = \widetilde{X} \cap \widetilde{Y}$
\[
  \begin{array}{l}
    X \underset{\mathcal{M}}{\between} Y \iff \forall\, X' \in \mathrm{omb}_{\Lambda}(X),\ Y' \in \mathrm{omb}_{\Lambda}(Y) \quad (X'_L = \emptyset,\ Y'_L = \emptyset) \Longrightarrow X' \underset{\Lambda}{\parallel} Y' \\
    X \underset{\mathcal{M}}{\parallel} Y \iff \forall\, X' \in \mathrm{omb}_{\Lambda}(X),\ Y' \in \mathrm{omb}_{\Lambda}(Y) ,\ \text{on a } X' \underset{\Lambda}{\parallel} Y' .
  \end{array}
\]
\marginal{\emph{Non}, cf.\ p.~49}
\note{note de sa main, « Non » doublement souligné, dans la marge gauche en face de « je dis que ça suffit » ; sous le signe de compatibilité de la première ligne, un indice noirci}

\page{48}
Ces relations sont symétriques, $\between$ est réflexive (tautologique), $\underset{\mathcal{M}}{\parallel}$ antiréflexive (utilisant le fait que $\underset{\Lambda}{\parallel}$ l'est, et M$\Lambda$2, impliquant que $\mathrm{omb}_{\Lambda}(X) \neq \emptyset$). Je dis

\textbf{\emph{Lemme 1}} Soient $X, Y \in \mathcal{M}$.
\begin{itemize}
  \item[a)] $X \underset{\mathcal{M}}{\parallel} Y \iff X \between Y$, et $\underbrace{\widetilde{X} \cap \widetilde{Y}}_{L} = \emptyset$
  \item[b)] Si $X, Y \in \Lambda$, alors $X \underset{\mathcal{M}}{\parallel} Y \iff X \underset{\Lambda}{\parallel} Y$ (on omettra donc par la suite les indices $\mathcal{M}$, $\Lambda$ dans $\parallel$)
\end{itemize}
\note{sous $\widetilde{X} \cap \widetilde{Y}$, un trait ondulé porte « $L$ »}

\emph{Preuve} a). \struck{On a $\Longrightarrow$ : $X \between Y$ trivialement par définition, $\widetilde{X} \cap \widetilde{Y} = \emptyset$}

Preuve $\Longrightarrow$ i.e.\ supposons $X \underset{\mathcal{M}}{\parallel} Y$, prouvons :
\begin{itemize}
  \item[1°)] $X \between Y$ : tautologique
  \item[2°)] $\widetilde{X} \cap \widetilde{Y} = \emptyset$. Si on avait $Z \in \widetilde{X} \cap \widetilde{Y}$, soit $Z' \in \mathrm{omb}_{\Lambda}(Z)$, alors $Z' \in \mathrm{omb}_{\Lambda}(X) \cap \mathrm{omb}_{\Lambda}(Y)$ \struck{donc que}. $X \parallel Y$ on doit avoir $Z' \parallel Z'$, absurde (antiréflexivité).
\end{itemize}

Preuve $\Longleftarrow$ i.e.\ supposons $X \between Y$, $\widetilde{X} \cap \widetilde{Y} = \emptyset$, et prouvons $X \parallel Y$. \struck{Soit $X' \in \mathrm{omb}_{\Lambda}(X)$, $Y' \in \mathrm{omb}_{\Lambda}(Y)$.} \add{Mais comme} \struck{Par def} \struck{\ill{} $X \between Y$, donc} \ldots\ on a $L = \widetilde{X} \cap \widetilde{Y} = \emptyset$, \struck{elle} donc $\mathrm{omb}_{\Lambda}(L) = \emptyset$, \struck{elle se réduit : $\forall\, X' \in \mathrm{omb}_{\Lambda}(X)$} la définition de $X \between Y$ montre qu'il se réduit : $X \parallel Y$.

\emph{Preuve} b) Résulte tautologiquement de M$\Lambda$4.

\emph{Terminologie} Soient $P, Q$ deux parties de $\mathcal{M}$. On pose
\[
  \begin{aligned}
    P \parallel Q &\overset{\mathrm{def}}{\iff} \forall\, X \in P,\ \forall\, Y \in Q,\ \text{on a } X \parallel Y \\
    &\iff \forall\, X' \in \mathrm{omb}_{\Lambda} P,\ Y' \in \mathrm{omb}_{\Lambda}(Q),\ \text{on a } X' \parallel Y'
  \end{aligned}
\]
(où $\mathrm{omb}_{\Lambda}(P) = \bigcup_{X \in P} \mathrm{omb}_{\Lambda}(X)$, et item pour $\mathrm{omb}_{\Lambda}(Q)$)

\page{49}
\textbf{\emph{Lemme 2}} Soient $X, Y \in \mathcal{M}$, avec $X \between Y$, $L = \widetilde{X} \cap \widetilde{Y}$, $X' \underline{\ll} X$, $Y' \underline{\ll} Y$. Posons $X'_L = \widetilde{X}' \cap \mathrm{Omb}(L)$, $Y'_L = \widetilde{Y}' \cap \mathrm{Omb}(L)$. Alors
\[
  X'_L \parallel Y'_L \iff X' \parallel Y'
\]
\note{les deux signes « $\ll$ » sont soulignés, comme page 41}

La partie \struck{\uncertain{tautologique}} \struck{\ill{}} \struck{$\Longrightarrow$} \struck{\uncertain{est}} $\Longrightarrow$ est bénigne. Je n'arrive pas à la prouver en utilisant seulement M$\Lambda$5 — il semble qu'il faut renforcer M$\Lambda$5 par cette forme plus forte, un peu décourageante par sa complication !
\note{au-dessus de la ligne biffée, trois mots interlinéaires, biffés eux aussi, non lus}

\uncertain{Ce} \uncertain{lemme} \uncertain{est} \uncertain{le} M$\Lambda$4 \uncertain{de} p.~47. Il faut \uncertain{prouver} \uncertain{ce} M$\Lambda$5 de la \uncertain{page} :

\textbf{\emph{Lemme 3}} Soient $X, Y$ avec $X \between Y$, $L = \widetilde{X} \cap \widetilde{Y}$.
\struck{$L = \widetilde{X} \cap \widetilde{Y}$, $Z \in \mathrm{omb}_{\Lambda}(X) \cap \mathrm{omb}_{\Lambda}(Y)$} \struck{$Z \in \Lambda$ tel que $Z \ll X$,}
Alors $\mathrm{omb}_{\Lambda}(X) \cap \mathrm{omb}_{\Lambda}(Y) = \mathrm{omb}_{\Lambda}(L)$.

Je n'y arrive pas non plus. Il faut donc \uncertain{ajouter} ces deux conditions, sur la relation $\underset{\Lambda}{\parallel}$ (en termes de la relation $\underset{\mathcal{M}}{\between}$ que $\parallel$ définit, cf.\ page précédente). Donc il faut mettre un M$\Lambda$6 en conséquence. Mais alors, en termes de la proposition p.~42, ça revient tautologiquement à dire qu'on a un système $(\mathcal{M}, \leq, \struck{\ll}, \between, \Lambda)$, \struck{satisfaisant} \struck{provenant d'un atelier qui satisfait $\Lambda 1$, $\Lambda 2$} \add{spécieux}. Rien de nouveau ! On \uncertain{tourne} \uncertain{en rond} !
\note{les trois dernières lignes sont serrées au bas du feuillet ; « spécieux » est écrit au-dessus de la partie biffée ; le symbole biffé dans le quadruplet est un $\ll$ (ou un $\between$) raturé}

\page{50}
\subsection{Récapitulation}
\note{titre de sa main, souligné d'un trait ondulé, en tête de la page 50}

Les axiomes algébriques d'ateliers, en termes de la donnée $(\mathfrak{F}, \leq, \ll)$, sont récapitulés p.~23. Disons un mot d'implication, concernant les axiomes d'incompatibilité-disjonction, voir p.~22 ou p.~45, 46 (avec une partie $\Lambda \subset \mathcal{M}$ dense : la clef, qui conduit à la notion d'atelier $\Lambda$-spécieux ou $\Lambda$-polyédral…). Dans le cas « spécieux » ou « polyédral », j'implique maintenant l'axiome At C (spéc) resp.\ At C (pol).
\note{« p.~23 », « p.~22 » renvoient à sa pagination de ce chapitre (lot 2), « p.~45, 46 » aux pages 45-46 ci-dessus}

\emph{Dans} le cas d'une structure $(\mathcal{M}, \leq, \ll, \between)$, les données sont plus maniables (quitte encore, j'ai supposé $\widetilde{\mathfrak{F}} \subset \mathfrak{P}(\mathcal{M})$ connue, \uncertain{donnée} compliquée !), mais le syst.\ d'axiomes pertinents semble plus simple, dans certains \uncertain{parmi} autres projets d'axiomes \emph{en} \uncertain{général}.
\note{« $\widetilde{\mathfrak{F}} \subset \mathfrak{P}(\mathcal{M})$ » est écrit au-dessus de la ligne et rattaché par un trait courbe}

Il y a lieu de considérer $(\mathcal{M}, \leq, \ll, \between)$ comme une structure algébrique indépendante, qu'on pourrait l'appeler un « pré-atelier » (ou « algèbre de multistrates »). Elle peut correspondre à \emph{plusieurs} ateliers (\uncertain{réalisés} comme des $\mathfrak{F} \subset \mathfrak{P}(\mathcal{M})$)

\page{51}
parmi lesquels il y en a un le plus grand de tous, formé de \emph{toutes} les parties \add{de $\mathcal{M}$} fermées pour $\leq$, et formées d'éléments deux à deux compatibles (pour $\between$).

J'distingue trois des conceptions possibles pour les axiomes.

\textbf{(1)} \add{Introduction} \struck{Axiomes} At 1 à At 7 sans plus, première version (cf.\ p.~6 \struck{\ill{}} — \uncertain{je} laisse tomber les axiomes relatifs à f)
\note{« (1) » est cerclé sur la page ; « Introduction » est écrit au-dessus de « Axiomes », biffé ; « p.~6 » est suivi d'un chiffre raturé. Ce renvoi est à sa pagination de ce chapitre (lot 1)}

\struck{At 11 1}
\begin{itemize}
  \item[\textbf{\emph{Prat 1}}] $X \leq Y \Rightarrow X \ll Y$
  \item[\textbf{\emph{Prat 2}}] $X \overset{\circ}{\ll} Y$ et $Y \overset{\circ}{\ll} Z \Rightarrow X \overset{\circ}{\ll} Z$
  \item[\textbf{\emph{Prat 3}}] \struck{Pour} $X \ll Y$, $\exists\, Y' \leq Y$ avec $X \overset{\circ}{\ll} Y'$, i.e.\ $X \in \mathrm{Omb}(Y')^{\circ}$ ; $\between$ réflexive, symétrique, et si $X' \leq X$, $Y' \leq Y$, $X \between Y \Rightarrow X' \between Y'$
  \item[\textbf{\emph{Prat 4}}] $X \overset{\circ}{\ll} Y$, $X \overset{\circ}{\ll} Y'$, $Y \between Y' \Longrightarrow Y = Y'$
  \item[\textbf{\emph{Prat 6}}] Si $X \between Y$, alors $X \ll Y \Longrightarrow X \leq Y$
\end{itemize}
\note{« Prat » (pré-atelier) est souligné à chaque fois ; les chiffres de « Prat 2 » et « Prat 3 » sont écrits en surcharge (l'un sur l'autre ?) ; « Prat 4 » est écrit au-dessus de « Prat 5 », qui n'est pas biffé nettement ; le « 6 » de « Prat 6 » est repassé. La ligne sur $\between$ n'a pas d'étiquette propre, elle est serrée sous Prat 3}
\marginal{\ill{} \ill{} transitivité des raff.\ intérieurs \ill{} — Axiome \ill{} \ill{} \ill{} — Unicité de la strate minimale \ill{} — \ill{} \ill{} $\ll$ \ill{} \ill{} $\leq$ \ill{} compatibles}
\note{notes obliques dans la marge gauche, en face de Prat 1-Prat 6, lues par fragments ; elles donnent le nom de chaque axiome}

NB seuls les \add{trois} \struck{deux} derniers axiomes font intervenir $\between$.

\textbf{(2)} Traduction At 1 à At 7 sans plus, mais en introduisant de plus une partie $\Lambda \subset \mathcal{M}$, dont on exige seulement qu'elle soit dense (i.e.\ $\forall\, X \in \mathcal{M}$, $\exists\, Z \in \Lambda$, $Z \overset{\circ}{\ll} X$, i.e.\ $\mathrm{omb}_{\Lambda}(X)^{\circ} \neq \emptyset$). (Cf.\ p.~42)
\note{« (2) » est cerclé}
\marginal{[NB. Le syst.\ d'axiomes \ill{} \ill{} \ill{} \ill{} \uncertain{précédent} \ill{} $\Lambda = \mathcal{M}$]}
\marginal{[NB 1, 3, 4 sont comme pour Prat (sauf \uncertain{somme} $\Lambda$ pour 3]}

\begin{itemize}
  \item[\textbf{\emph{Prat$\Lambda$1}}] $X \leq Y \Rightarrow X \ll Y$
  \item[\textbf{\emph{Prat$\Lambda$2}}] $\forall\, X \in \mathcal{M}$, $\exists\, Z \in \Lambda$ avec $Z \overset{\circ}{\ll} X$ (i.e.\ $\mathrm{omb}_{\Lambda}(X)^{\circ} \neq \emptyset$)
  \item[\textbf{\emph{Prat$\Lambda$3}}] $\forall\, X \in \mathcal{M}$, $Z \in \struck{\ill{}}\, \Lambda$, $Z \ll X$, $\exists\, X' \leq X$, avec $Z \overset{\circ}{\ll} X'$
  \item[\textbf{\emph{Prat$\Lambda$5}}] $\between$ réflexive et symétrique, \add{si} $X' \leq X$, $Y' \leq Y$, \add{alors} $X \between Y \Rightarrow X' \between Y'$
  \item[\textbf{\emph{Prat$\Lambda$6}}] Si $X \between Y$, et $Z \in \Lambda$ telle que $Z \ll X$, $Z \ll Y$, alors $\exists\, Z' \in \mathcal{M}$ avec $Z' \ll Z \leq X, Y$ [i.e.\ $\mathrm{omb}_{\Lambda}(X) \cap \mathrm{omb}_{\Lambda}(Y) = \mathrm{omb}_{\Lambda}(L)$, où $L = \widetilde{X} \cap \widetilde{Y}$]
\end{itemize}
\note{les étiquettes « Prat$\Lambda$5 » et « Prat$\Lambda$6 » sont écrites en surcharge (sur « 4 » et « 5 » ?) ; une accolade à gauche, partant de Prat$\Lambda$4 plus bas, remonte par une flèche jusqu'à Prat$\Lambda$5 : l'axiome Prat$\Lambda$4 est à insérer avant. Dans Prat$\Lambda$6, la page écrit « $Z' \ll Z \leq X, Y$ », alors que M$\Lambda$5 (p.~42) a $Z \ll Z' \leq X, Y$}
\marginal{(évident si $\Lambda = \mathcal{M}$) — strate minimale pour les $Z \in \Lambda$}
\marginal{Axiome des raff.\ induits}
\marginal{On peut \uncertain{bloquer} Pr$\Lambda$3, $\Lambda$4 en \ill{} \ill{} $\forall\, Z \ll X$ \ill{} \ill{} \ill{} \ill{} \ill{}}
\note{notes obliques dans la marge gauche, en face de Prat$\Lambda$2-Prat$\Lambda$4, lues par fragments}

Seuls les deux derniers axiomes font intervenir $\between$.

\begin{itemize}
  \item[\textbf{\emph{Prat$\Lambda$4}}] Soient $Y \underline{\ll} X$, $X' \ll X$, et $Z \in \Lambda$ telle que $Z \ll Y, X'$, alors $\exists\, Y'$ avec $Z \ll Y'$, $Y' \leq Y$, $Y' \underline{\ll} X'$. Si $Y \leq X$, alors $Y' \leq X'$. [i.e.\ $\mathrm{omb}_{\Lambda}(Y) \cap \mathrm{omb}_{\Lambda}(X') = \mathrm{omb}_{\Lambda}(Y_{X'})$ où $Y_{X'} \overset{\mathrm{def}}{=} \widetilde{Y} \cap \widehat{\mathrm{Omb}}(X')$]
\end{itemize}
\note{la relation entre $X'$ et $X$ est un $\ll$ vertical, $X'$ sous $X$, comme page 42 ; les deux $\ll$ soulignés comme page 41 ; « i.e. » est entouré ; la dernière formule, encadrée, est serrée au bas du feuillet et le chapeau sur « Omb » n'est pas sûr}

\page{52}
\note{son numéro « 52 » est repassé et entouré}
\textbf{(3)} Cas des ateliers spécieux ou polyédraux, axiomatisés en termes de $(\mathcal{M}, \leq, \ll, \underset{\mathcal{L}}{\parallel})$, où $\underset{\mathcal{L}}{\parallel}$ est une relation « de disjonction » sur $\mathcal{L}$ = ens.\ des él.\ min.\ de $(\mathcal{M}, \ll)$ (cf.\ p.~34)
\note{« (3) » est cerclé ; « p.~34 » renvoie à sa pagination de ce chapitre (lot 2)}
\marginal{[Prat dis 1 à Prat dis 5 comme Prat$\Lambda$1 à Prat$\Lambda$5, avec $\Lambda = \mathcal{L}$]}
\marginal{« dis » pour « disjonction » — « Prat » pour « préatelier »}
\note{la première note est entre crochets à droite du titre ; la seconde, oblique, dans la marge gauche, lue en partie}

\begin{itemize}
  \item[\textbf{\emph{Prat dis 1'}}] $X \leq Y \Longrightarrow X \ll Y$
  \item[\textbf{\emph{Prat dis 2'}}] \struck{Comme} $\forall\, X \in \mathcal{M}$, $\mathrm{omb}(X)^{\circ} \neq \emptyset$ i.e.\ $\exists\, x \in \mathcal{L}$, $x \overset{\circ}{\ll} X$
  \item[\textbf{\emph{Prat dis 3'}}] $\forall\, X \in \mathcal{M}$, $x \in \mathcal{L}$, $x \ll X$, $\exists\, X' \in \mathcal{M}$, $x \overset{\circ}{\ll} X' \leq X$
  \item[\textbf{\emph{Prat dis 4'}}] Comme Prat$\Lambda$4 : i.e.\ Soit $Y \ll X$, $X' \ll X$ et $x \in \mathcal{L}$, $x \ll Y, X'$, alors $\exists\, Y'$ avec $Y' \leq Y$, $Y' \ll X'$, $x \ll Y'$, i.e.
  \[
    \mathrm{omb}(Y) \cap \mathrm{omb}(X') = \mathrm{omb}(Y_{X'}) \overset{\mathrm{(def)}}{=} \bigcup_{Z \in Y_{X'}} \mathrm{omb}(Z) \quad \text{où } Y_{X'} = \widetilde{Y} \cap \mathrm{Omb}(X') .
  \]
  De plus, si $Y \leq X$, alors $Y' \leq X'$, i.e.\ $\mathrm{omb}(Y) \cap \mathrm{omb}(X') = \mathrm{omb}(L)$ où $L = \widetilde{Y} \cap \widetilde{X}'$.
  \item[\textbf{\emph{Prat dis 5'}}] $\underset{\mathcal{L}}{\parallel}$ est \struck{symétrique, réflexive} \add{symétrique} et antiréflexive (i.e.\ $x \parallel y$ implique $y \parallel x$ et $x \neq y$).
  \item[\textbf{\emph{Prat dis 6'}}] $\forall\, X, Y \in \mathcal{M}$ avec $Y \leq X$, \struck{\ill{}} $x, y \in \mathcal{L}$, $y \ll Y$, $x \ll X$ et \struck{$y \in \mathrm{omb}$} $x \not\ll Y$ (i.e.\ $y \in \mathrm{omb}(Y)$, $x \in \mathrm{omb}(X) \smallsetminus \mathrm{omb}(Y)$) on a $x \parallel y$.
\end{itemize}
\note{dans chaque étiquette de 2' à 6', « dis » est biffé et surmonté d'un « $\mathcal{L}$ » ; dans 1', le mot après « Prat » est noirci. Dans Prat dis 4', $X'$ est écrit sous $X$ avec un $\ll$ vertical, comme page 42}
\marginal{[Seules les deux dernières axiomes font intervenir $\parallel$]}

Ce sont les mêmes axiomes que Prat$\Lambda$, à l'exception du dernier Prat dis 6 (et à l'exception du fait que Prat dis 5 prend une forme plus simple) — et qui remplace les derniers axiomes en termes d'incompatibilité.
\note{« en termes d'incompatibilité » : lecture incertaine du dernier mot}

On définira :
\[
  X \between Y \iff
  \begin{cases}
    \forall\, x \in \mathrm{omb}(X) \smallsetminus \mathrm{omb}(L),\ y \in \mathrm{omb}(Y) \smallsetminus \mathrm{omb}(L) \ (\text{où } L = \widetilde{X} \cap \widetilde{Y}) \text{ on a } x \parallel y \\
    [+\ L = \emptyset \text{ ou } L \text{ de la forme } \widetilde{Z}, \text{ dans le cas « polyédral »}]
  \end{cases}
\]
\marginal{Donc le \emph{vrai} \uncertain{problème} d'axiomatisation pour \uncertain{un} atelier (de plus) est de \uncertain{construire} \ill{} \ill{} \ill{} \ill{} \ill{} \ill{}}
\note{note oblique dans la marge inférieure gauche, « vrai » triplement souligné, lue par fragments ; la page est barrée en bas à droite d'un trait oblique}

\page{53}
Mais dans le cas polyédral, il faut un axiome suppl.\ :

\textbf{\emph{Prat pol}} $\forall\, X \in \mathcal{M}$, $\widetilde{X}$ est polyédral, i.e.\ si $Y, Z \in \mathcal{M}$ et si $\lbrace X, Y \rbrace$ est majoré et minoré dans $\mathcal{M}$, alors $\mathrm{Inf}^{\leq}(Y, Z)$ existe (i.e.\ $\widetilde{X} \cap \widetilde{Y}$ a un plus grand élément).
\note{l'énoncé est marqué d'un trait vertical dans la marge ; la page écrit « $\lbrace X, Y \rbrace$ » puis « $\mathrm{Inf}^{\leq}(Y, Z)$ », où l'on attend les mêmes lettres ; dans $\mathrm{Inf}^{\leq}(Y, Z)$ les deux lettres sont repassées}

\emph{NB} Alors $\struck{\mathrm{Sup}}\ \mathrm{Inf}^{\leq}(Y, Z)$ existe aussi \emph{dans} $\widetilde{X}$ (mais \emph{pas} \uncertain{nécessairement} dans $\mathcal{M}$…)

\emph{Remarque} Soit $(\mathcal{M}, \leq)$ un ens.\ ordonné, et prenons $(\ll) \Leftrightarrow (\leq)$, $(x \underset{\mathcal{L}}{\parallel} y) \Leftrightarrow (x \neq y)$. Alors tous les axiomes sont satisfaits, sauf le seul axiome de divisibilité Prat$\mathcal{L}$2'. Je pense p.\ ex.\ au cas où $\mathcal{M}$ = ensemble des simplexes d'un ens.\ quasi simplicial, $(\mathcal{L}, \mathcal{M} \subset \mathfrak{P}_{\mathrm{f}}(\mathcal{L}))$, où les figures sont les sous-ens.\ quasi simpliciaux, correspondant aux sous-ensembles fermés de $\lbrace \mathcal{M}, \leq (\overset{\mathrm{def}}{=} \subset) \rbrace$. Et en termes de l'axiomatique (1) ou (2), en prenant
\[
  X \ll Y \iff X \leq Y , \quad \text{et } X \between Y \text{ tjrs satisfaits.}
\]
\struck{$X \overset{\circ}{\ll} Y$ signifie $X = Y$.} Appelons « trivial » un préatelier tel que $X \ll Y \iff X \leq Y$, $X \between Y$ tjrs satisfaits, de sorte que
\[
  X \overset{\circ}{\ll} Y \iff X = Y
\]
\note{« Prat$\mathcal{L}$2' » : $\mathcal{L}$ pour « dis », comme page 52 ; « un préatelier » est écrit sous un mot interlinéaire (« imaginaire » ?) non lu}

Alors tous les axiomes Prat 1 - Prat 6 sont automatiquement satisfaits — on peut donc partir d'un ens.\ ordonné $(\mathcal{M}, \leq)$ quelconque. Passant à $\mathfrak{F}$, on trouve \struck{\ill{}}
\[
  \mathfrak{F} = \mathfrak{P}_{\mathrm{f}}(\mathcal{M}, \leq) \quad \text{ensemble de \emph{toutes} les parties fermées}
\]
dans $\mathfrak{F}$, $F \ll G \iff F \leq G$, et $F \between G$ tjrs satisfait, $F \parallel G \iff F \cap G = \emptyset$. \struck{Axiomes C et} Les conditions At D, At C (spéc), At Fil, At supp sont satisfaites. Pour les conditions sur les lieux, At L 2 \struck{At L 1} qui est équivalent : Tout \struck{\ill{}} $X \in \mathcal{M}$ est minoré par un élément minimal — n'est pas prédisant ! L'axiome At L 3 tjrs satisfait (car deux lieux distincts sont disjoints, itou compatibles). Donc le seul axiome pas satisfait est At L 1, de divisibilité.
\note{« At Fil » : lecture incertaine du dernier mot ; « prédisant » : lecture incertaine. Les trois dernières lignes sont serrées au bas du feuillet}
\marginal{On trouve \ill{} \ill{} \ill{} \ill{} \ill{} \ill{} \ill{} \ill{} \ill{} \ill{} \ill{} — NB \ill{} « \ill{} \ill{} » \ill{} \ill{} \ill{} !}
\note{longue note oblique dans la marge inférieure gauche, couverte de biffures et de traits ; illisible sauf quelques mots}

\page{54}
\subsection{25 juin}
\marginal{25 juin}
\note{date de sa main, écrite en oblique dans l'angle supérieur gauche de la page 54, soulignée d'un trait qui la barre en partie, en face du premier alinéa}

Je voudrais revenir sur la notion d'\emph{atelier} \emph{ensembliste}. Ce sera un atelier soit spécieux, soit polyédral, prenons le premier cas.
\note{il reprend ici les ateliers ensemblistes de la deuxième mouture (GF VI, pp.~52-58 : At ens 1, Atens*1-3, la Scholie de la p.~54), désormais caractérisés par le morphisme $\varphi$ vers $\mathrm{Fig}(\mathcal{L})$}

\textbf{\emph{Proposition}} Soit $\mathcal{A} = (\mathfrak{F}, \leq, \ll)$ un atelier, considérons l'(ho)morphisme canonique
\[
  \varphi : \mathfrak{F} \longrightarrow \mathrm{Fig}(\mathcal{L}) \qquad F \longmapsto \mathrm{multomb}(F)
\]
\add{Considérons les} \struck{Conditions : équivalentes}
\begin{itemize}
  \item[(i)] L'hom.\ $\varphi$ est un \emph{plongement} d'ateliers (cf.\ GF VI p.~51) et l'image de $\mathfrak{F}$ par $\varphi$ est un sous-atelier \emph{spécieux} (\add{p.~40}) (localité : si $F, G \in \varphi(\mathfrak{F})$ et $F \between G$ dans $\mathrm{Fig}(\mathcal{L})$, alors $F \between G$ dans $\varphi(\mathfrak{F})$ i.e.\ $F \cup G \in \varphi(\mathfrak{F})$).
\end{itemize}

En d'autres termes :
\begin{itemize}
    \item[a)] $\forall\, F \in \mathfrak{F}$, $F' \mapsto \varphi(F')$ induit une bijection $\mathrm{Ssfig}(F) \longrightarrow \mathrm{Ssfig}(\varphi(F))$
    \item[b)] $\varphi(F) \ll \varphi(G) \Longrightarrow F \ll G$ \quad (donc $\varphi$ injective)
    \item[c)] Si $\varphi(F) \between \varphi(G)$, $\Longrightarrow F \between G$
\end{itemize}
\begin{itemize}
  \item[(ii)] $\mathfrak{F}$ est \emph{spécieux} (au sens absolu) i.e.\ satisfait \add{les} (pour $\Lambda = \mathcal{L}$) quatre systèmes équivalents de conditions en la p.~45 (p.\ ex.\ At L1, At L2, At L3, At CL (spéc)), \emph{et} pour $x, y \in \mathcal{L}$, $x \neq y \Longrightarrow x \overline{\parallel}\, y$ (donc At L 3 devient automatique).
\end{itemize}
Alors (i) $\Longrightarrow$ (ii) \add{(ii)} $\Longrightarrow$ (i a) \struck{\ill{}} (mais pas nécessairement (i b), vu l'injectivité de $\varphi$)
\note{« Mulomb(F) » sur la page, d'une plume repassée, lu $\mathrm{multomb}(F)$ ; « GF VI p.~51 » renvoie à la deuxième mouture (folder 156-6), page 51, où sont définis les plongements d'ateliers (conditions a'), b')) ; « p.~40 » est écrit au-dessus de « spécieux », et les pages 38-40 de GF VI définissent sous-ateliers et parties « spécieuses » ; « p.~45 » renvoie à la page 45 ci-dessus. Le signe après $x \neq y$ est un $\parallel$ barré de traits horizontaux, lu $\overline{\parallel}$ (non disjoints) ; « (ii) » sous « $\Longrightarrow$ (i a) » est un ajout interlinéaire ; au-dessus de « (i b) », « \ill{} (i, c) » est écrit et rattaché par une flèche}
\marginal{$\varphi$ plongement}
\marginal{$\varphi$ plongement « spécieux »}
\note{deux notes obliques dans la marge gauche, en face de a)-b) et de c)}

\emph{Dém} (i) $\Longrightarrow$ (ii). \struck{Comme} Comme on sait que $x, y \in \mathcal{L}$ \struck{$(x, y) \in \mathcal{L}$} $\Longrightarrow x \between y$, $x, y \in \mathcal{L}$, $x \neq y \Longrightarrow \struck{\ill{}}\ x \parallel y$ \ldots\ Comme $\varphi(x)$, $\varphi(y)$ sont des \struck{\ill{}} figures ponctuelles i.e.\ $\varphi(x) \between \varphi(y)$ donc $x \between y$ par c). At L1 : Soit $X \in \mathcal{M}$, \struck{soit $x$} comme $X \neq \emptyset$ i.e.\ $\varphi(X) \neq \emptyset$, donc par définition de $\varphi$ (ceci signifie $\mathrm{omb}(X)^{\circ} \neq \emptyset$), OK.
\note{ces lignes de démonstration sont très abrégées ; la première implication « $x, y \in \mathcal{L}$, $x \neq y \Longrightarrow x \parallel y$ » est celle qu'on lit, sans que le raccord avec ce qui suit soit sûr}

\page{55}
At L 2 : Soient $X, Y \in \mathcal{M}$, avec $\forall\, x \in \mathrm{omb}(X)$, $y \in \mathrm{omb}(Y)$ on a $x \neq y$ \add{($\Longrightarrow x \parallel y$)} i.e.\ $\mathrm{omb}(X) \cap \mathrm{omb}(Y) = \emptyset$. \struck{Ceci signifie} \struck{que $\mathrm{omb}(X) \cap \varphi(X) \parallel$} Comme $\mathrm{omb}(X) = |\varphi(X)|$ et item pour $\mathrm{omb}(Y)$, ceci exprime donc $\varphi(X) \parallel \varphi(Y)$. Injectivité \struck{\ill{}} $\varphi(X) \between \varphi(Y)$ donc $X \between Y$. Il est évident d'autre part que $X \cap Y = \emptyset$ (à cause de L 1 ; si $Z \leq X, Y$, $\exists\, x \in \mathcal{L}$, $x \ll Z$, donc $x \in \mathrm{omb}(X) \cap \mathrm{omb}(Y)$, absurde).
\note{« ($\Longrightarrow x \parallel y$) » est écrit au-dessus de « $x \neq y$ » ; la ligne biffée est recouverte d'un trait. La phrase « Injectivité … donc $X \between Y$ » se lit ainsi ; on attend « par c) », comme page 54}

At CL (spéc) : Soient $X, Y \in \mathcal{M}$, $L = \widetilde{X} \cap \widetilde{Y}$, supposons que $\forall\, x \in \mathrm{omb}(X) \smallsetminus \mathrm{omb}(L)$, $y \in \mathrm{omb}(Y) \smallsetminus \mathrm{omb}(L)$, on ait $x \parallel y$ i.e.\ $x \neq y$, on déduit en termes \uncertain{ensemblistes}
\[
  \mathrm{omb}(X) \cap \mathrm{omb}(Y) = \struck{\emptyset}\ \mathrm{omb}(L)
\]
Prouvons $X \between Y$. On \struck{\ill{}} \emph{sait} \uncertain{que} la condition est
\[
  |\varphi(X)| \cap |\varphi(Y)| = |\varphi(L)|
\]
où $\varphi(L)$ est une sous-fig.\ ensembliste commune de $\varphi(X)$, $\varphi(Y)$, donc $\varphi(X) \between \varphi(Y)$, donc $X \between Y$, OK.
\note{« sait » est souligné, à la place d'un mot biffé}

(ii) $\Longrightarrow$ (i) \add{a, c)} Il faut prouver a) \struck{\ill{}} c)
\note{« a, c) » est écrit au-dessus de « (i) » et rattaché par un trait courbe}

\textbf{\emph{Lemme 1}} Soit $\mathcal{A} = (\mathfrak{F}, \leq, \ll)$, \add{atelier, $\Lambda$ partie dense de $\mathcal{M}$}, \struck{satisfaisant L 1. Alors}, \struck{\ill{}} Considérons le \struck{\ill{}} morphisme canonique
\[
  \varphi_{\Lambda} : \mathfrak{F} \longrightarrow \mathrm{Fig}(\Lambda) \qquad F \longmapsto \mathrm{multomb}_{\Lambda}(F)
\]
Alors $\varphi_{\Lambda}$ satisfait a), i.e.\ $\forall\, F$, on a
\[
  \mathrm{Ssfig}(F) \xrightarrow{\ \sim\ } \mathrm{Ssfig}(\varphi_{\Lambda}(F)) , \quad F' \longmapsto \varphi_{\Lambda}(F')
\]
Je \uncertain{voudrais} en \uncertain{effet} \uncertain{dire} que l'application $\widetilde{F} \longrightarrow \widetilde{\varphi_{\Lambda}(F)}$ est un iso.\ d'ens.\ ordonnés
\marginal{La densité de $\Lambda$ sert à assurer que $\varphi_{\Lambda}$ \ill{} des strates \ill{} \ill{} \ill{} strates}
\note{note oblique dans la marge inférieure gauche, en face de « Alors $\varphi_{\Lambda}$ satisfait a) », avec une flèche vers le haut ; lue en partie. Le tilde de $\widetilde{\varphi_{\Lambda}(F)}$ est un trait courbe au-dessus de l'expression, surmonté d'un mot (« de » ?)}

\page{56}
Ça a été vu pour $\Lambda = \mathcal{M}$. Dans le cas général, c'est un fait ens., sa preuve paraît aussi aisée.
\[
  i : \Lambda \longrightarrow \mathcal{M}
\]
et figure ens.\ $\Phi \in \mathrm{Fig}(\mathcal{M})$, \uncertain{pour} $\forall\, X \in \Phi$, on a $i^{*}(X^{\circ}) \neq \emptyset$.
\note{la lettre bouclée pour une figure ensembliste est lue $\Phi$, comme dans GF VI ; cette première ligne, « Dans le cas général, c'est un fait ens., sa preuve paraît aussi aisée », est d'une lecture en partie incertaine}

Prenant $\Lambda = \mathcal{L}$ dans \struck{\ill{}} le cas de la prop.\ \add{(supposant (ii))} on trouve a.
\note{« (supposant (ii)) » est écrit au-dessus de la ligne et rattaché par un trait courbe}

Prouvons c). \struck{\ill{}} Supposons $\varphi(F) \between \varphi(G)$, prouvons $F \between G$. Compte tenu de a) il s'ensuit que la plus grande sous-figure commune de $\varphi(F)$, $\varphi(G)$ est $\varphi(L)$, où $L = F \cap G$ \struck{\ill{}} ($\overset{\mathrm{def}}{=} \mathrm{Inf}^{\leq}(F, G)$, donc $\widetilde{L} = \widetilde{F} \cap \widetilde{G}$). En effet, $\varphi(L)$ est une telle \struck{sous-}figure commune, et si $\Phi$ en est \struck{une} \add{la} \emph{plus grande} \struck{\ill{}} commune de $\varphi(F)$, $\varphi(G)$, on a $\varphi(L) \subset \Phi$, \struck{donc $\Phi$} et par a) on a $\Phi = \varphi(L') = \varphi(L'')$, où $L' \leq F$, $L'' \leq G$. \struck{Or} $L \leq L', L''$. Prouvons que $L' = L''$ donc $L' = L'' = L$. (\emph{NB} on ne sait pas si $\varphi$ est injectif, et $\varphi$ pourrait ne pas l'être…) \add{Si} c'était faux, soit $X$ \struck{\ill{}} $\in \struck{\ill{}}\ \widetilde{L}' \smallsetminus \widetilde{L}$ il existe (à cause de a)) un unique $Y \in \widetilde{L}'' \smallsetminus \widetilde{L}$ telles que $\varphi(X) = \varphi(Y)$. Alors la plus grande sous-figure commune de $X$, $Y$ est $X \cap L = Y \cap L$, distincte de $X$, $Y$. Est-ce bien impossible ? Si on trouve de telles situations, on aurait $\varphi(X) \between \varphi(Y)$ puisque $\varphi(X) = \varphi(Y)$, mais $X \between Y$ à cause de a) appliqué à \struck{\ill{}} $F = X \cup Y$ (si on avait $X \between Y$).
\note{« $X \in \widetilde{L}' \smallsetminus \widetilde{L}$ » : un mot biffé après $X$, puis un symbole noirci ; « $\in \widetilde{L}'' \smallsetminus \widetilde{L}$ » est écrit au-dessus de « telles que », à la place d'un $Y$ noirci. La conclusion « mais $X \between Y$ » se lit ainsi sur la page, alors que l'argument attend « $X$ et $Y$ non compatibles » ; la parenthèse finale « (si on avait $X \between Y$) » est de lecture incertaine}

\page{57}
Notons qu'on \struck{\ill{}} devrait avoir \struck{$\mathrm{omb}(X) \cap \mathrm{omb}(Y) \neq \emptyset$}
\[
  \mathrm{omb}(X) = \mathrm{omb}(Y) \quad (= |\varphi(X)| = |\varphi(Y)|)
\]
et qu'a priori il n'y a pas à exclure que $L = \emptyset$.

\emph{Exemple} $X$ espace topologique. On prend \struck{\ill{}}
\[
  \mathcal{M} = \text{simplexes \emph{affines} dans } X
\]
avec la relation $\leq$ habituelle, \struck{$\ll$ la relation} (mais en comparant les structures affines), $\ll$ à l'inverse. Nous exigeons que $|\sigma| \cap |\sigma'|$ soit vide ou un sous-simplexe commun (\ill{} \ill{}), et $\between$ itou (\ill{} \ill{} \ill{}), disons $\mathfrak{F}$, satisfaisant toutes les conditions. Mais considérons, sur une partie de $X$ homéomorphe à $[0, 1]$, deux structures affines différentes, disons $X, Y \in \mathcal{M}$, \struck{\ill{}} avec ici $L = \partial X = \partial Y$, $\varphi(X) = \varphi(Y)$, mais $X \neq Y$. Donc (ici $\varphi(X) \between \varphi(Y)$ et $\varphi(X) \ll \varphi(Y)$) \uncertain{(i a, c)} \uncertain{non} : ni $X \between Y$ ni $X \ll Y$.
\note{la page est de lecture difficile ici. « (\ill{} \ill{}) » : deux mots non lus au bout de la ligne ; « $\between$ itou » est suivi d'une parenthèse de trois mots non lus ; « (ii) » est écrit au-dessus de « conditions » et rattaché par un trait. La fin de la dernière phrase est serrée à gauche en bas du paragraphe ; « (i a, c) » est écrit en petit, en oblique ; le premier « ni » n'est pas sûr, la page donnant « $X \between Y$ » sans négation visible}

\textbf{\emph{Cor}} La condition (i) implique (ii) \add{+ l'injectivité de $\varphi$}. \struck{Mais réciproquement (ii) + l'injectivité de $\varphi | \mathcal{M}$} \struck{équivaut} pour \uncertain{mai}\ldots
\note{l'énoncé est marqué d'un trait vertical dans la marge ; « implique » remplace un mot biffé (« équivaut » ?) ; la ligne suivante est biffée ; « pour mai… » est inachevé}
\marginal{Plus précisément (i) $\Rightarrow$ (ii) + $\varphi$ injectif $\Rightarrow$ (i a), c)}
\note{note oblique dans la marge gauche, en face du Cor}

Bien sûr, (i) implique l'injectivité de $\varphi$, (ii) + $\varphi | \mathcal{M}$ inj.\ $\Longrightarrow$ b, c. Reste à prouver que \ldots\ (ça a été prouvé déjà).

Prouvons c), plutôt, en continuant. On a (utilisant l'injectivité de $\varphi$) que

\page{58}
$\varphi(L)$ est la plus grande sous-figure commune de $\varphi(F)$, $\varphi(G)$. L'hyp.\ $\varphi(F) \between \varphi(G)$ équivaut à :
\[
  |\varphi(F)| \cap |\varphi(G)| = |\varphi(L)| \quad \text{i.e.} \quad \mathrm{omb}(F) \cap \mathrm{omb}(G) = \mathrm{omb}(L) ,
\]
mais c'est (en vertu de At CL (spéc)) le condition pour que $F \between G$. OK.

Prouvons b), i.e.\ soient $F, G$ avec $\varphi(F) \ll \varphi(G)$, prouvons $F \ll G$. \struck{\ill{}} Utilisant a) déjà prouvé, on est ramené au cas $X, Y \in \mathcal{M}$. Mais c'est un \uncertain{faux} !
\note{« un faux ! » : lecture incertaine du mot avant le point d'exclamation, qui pourrait être « faux »}

\struck{\emph{Ex} $X$ espace affine sur $\mathbb{R}$. On prend $\mathcal{M}$ = simplexes affines, $\leq$ \ill{} \ill{}} \add{\uncertain{Cas} polyédral \ill{} comme dans l'exemple} \struck{précédent}, mais en remplaçant cette fois la relation $\ll$, par la relation ordinaire (« simplicial » — « polyédral »).
\note{ce paragraphe est barré de trois longs traits obliques ; la deuxième ligne biffée se termine par « sous-simplexes » ; l'ajout interlinéaire, au-dessus de « $\mathbb{R}$ », n'est lu qu'en partie}

Soit $\mathcal{L}$ un ens., considérons l'ens.\ $\mathrm{Figpol}(\mathcal{L})$ des figures \emph{polyédrales} dans $\mathcal{L}$. On y prend $\leq$ et $\ll$ comme d'habitude, \add{\uncertain{« ensembliste polyédral »}} étant entendu que $\ll$ signifie $\underset{\mathrm{pol}}{\ll}$, tandis qu'on prend la compatibilité $\between$ \emph{ensembliste} (pour son ensembliste non polyédral). (D'ailleurs, en fait, la $\ll$ ne modifie pas la validité des conditions Prat 1' -- Prat 6'.)
\note{« ensembliste polyédral » (lecture incertaine) est écrit au-dessus de « comme d'habitude » ; sous le deuxième $\ll$ de la phrase suivante est écrit « pol ». « ensembliste », après la compatibilité, est souligné}

Considérons deux \struck{figures} \add{strates} polyédrales (p.\ ex.\ en hyp.\ $\Delta_1$, $\Delta_n$) $X, Y$, \struck{\ill{}} avec $X \underset{\mathrm{ens}}{\ll} Y$ et \struck{\ill{}} \uncertain{hyp} non $X \underset{\mathrm{pol}}{\ll} Y$, on a\ldots\ Voilà le contre-exemple.
\note{« strates » est écrit au-dessus de « figures », biffé ; l'indice « ens » sous le premier $\ll$ de la dernière phrase est une lecture incertaine d'un mot court ; « hyp » suit deux mots biffés}

\page{59}
\textbf{\emph{Déf}} Un atelier est dit « ensembliste spécieux » s'il satisfait la condition (i). Idem pour « ensembliste polyédral », en remplaçant $\mathrm{Fig}(\mathcal{L})$ par $\mathrm{Figpol}(\mathcal{L})$. On a les \struck{conditions} notions \add{correspondantes} pour les préateliers.

Un atelier est dit « quasi-ensembliste » s'il satisfait (ii). Notion \add{analogue} pour le cas « quasi-ensembliste polyédral », ou le cas des préateliers.
\note{« correspondantes » et « analogue » : deux mots interlinéaires, lus d'après le sens}
\marginal{NB. Il faudrait encore \ill{} \ill{} \ill{} \ill{} \ill{} des figures dans \ill{} \ill{} sous-ensemblistes (\ill{} de figures ens.\ qui ne soient pas \ill{} \ill{} des figures \ill{} \ill{} \uncertain{recouvrements} \ill{})}
\note{longue note oblique dans la marge gauche, en face de la Déf., lue par fragments}

Les préateliers quasi-ensemblistes (i.e.\ affirmement spécieux ou polyédraux) correspondent aux $(\mathcal{M}, \leq, \ll)$ satisfaisant les conditions suivantes (déduites de Prat 1' -- Prat 6') :
\note{« affirmement » est sur la page, lu tel quel}

\begin{itemize}
  \item[\textbf{\emph{Prat quens 1}}] $X \leq Y \Longrightarrow X \ll Y$
  \item[\textbf{\emph{Prat quens 2}}] $\forall\, X \in \mathcal{M}$, $\mathrm{omb}(X)^{\circ} \neq \emptyset$ i.e.\ $\exists\, x \in \mathcal{L}$, $x \overset{\circ}{\ll} X$
  \item[\textbf{\emph{Prat quens 3}}] $\forall\, X \in \mathcal{M}$, $x \in \mathcal{L}$, $x \ll X$, $\exists\, X' \in \mathcal{M}$, $x \overset{\circ}{\ll} X' \leq X$
  \item[\textbf{\emph{Prat quens 4}}] Soit $Y \ll X$, $X' \ll X$, \struck{Alors} et $Y_{X'} = \widetilde{Y} \cap X'$. Alors
  \[
    \mathrm{omb}(Y) \cap \mathrm{omb}(X') = \mathrm{omb}(Y') \Bigl( \overset{\mathrm{def}}{=} \bigcup_{Z \in Y_{X'}} \mathrm{omb}(Z) \Bigr) ,
  \]
  i.e.\ si $x \in \mathcal{L}$, $x \ll Y, X'$, $\exists\, Y'$ avec $Y' \leq Y$, $Y' \ll X'$, $x \ll Y'$. De plus, si $Y \leq X$, alors $Y' \leq X'$ (i.e.\ on peut remplacer $Y_{X'}$ par $\widetilde{Y} \cap \widetilde{X}'$ — les deux devenant d'ailleurs égaux).
\end{itemize}
\note{« quens » (quasi-ensembliste) est souligné dans chaque étiquette. Dans Prat quens 4, $X'$ est écrit sous $X$, relié par un $\ll$ vertical, comme page 42 ; « $\mathrm{omb}(Y')$ » est écrit ainsi, là où la parenthèse définit $\mathrm{omb}(Y_{X'})$ ; « De plus, si » est suivi d'un mot biffé, puis « $Y' \leq X'$ » précédé d'un « OPS » (?) non sûr}
\marginal{densité divisibilité}
\marginal{strates minimales}
\marginal{raffinements induits}
\note{notes obliques dans la marge gauche, en face de Prat quens 2, 3, 4, reliées par des flèches}

Je n'ai plus besoin en effet d'introduire $\underset{\mathcal{L}}{\parallel}$, vu qu'on la connaît : $(x \parallel y) \iff (x \neq y)$.

\page{60}
Pratiquement, quand on ne travaille pas avec des « lieux » proprement dits (et je \uncertain{suis} le premier, \uncertain{visiblement}, à y \ill{} \ill{}, aux lieux\ldots) mais avec des points, et dans l'optique des espaces « infiniment divisibles » (et \ill{} \ill{}, je suis le premier \ill{} \ill{} \ill{}), il semblerait que cette structure \uncertain{satisfasse} à \ill{}, \uncertain{bien} \uncertain{sûr}, des besoins. \struck{\ill{} \ill{}} Cette structure a l'avantage d'être \emph{commune} au contexte « spécieux » et « polyédral », et \ill{} \emph{plus}, \ill{} \ill{} \ill{} (seulement $\leq$ et $\ll$ \ill{} \ill{} \ill{} \ill{} dans le cas des ateliers $\mathfrak{F}$) axiomes, réduits ici au nb de quatre \ill{} \ill{} plus compliqués \ill{} \ill{}. (Quand \ill{} travaille avec \ill{} \ill{} \ill{} \ill{} \ill{}, \ill{} \ill{} figures \ill{} \ill{} \ill{} \ill{} préatelier correspondant.) \ill{} \ill{} \ill{} \ill{} \ill{} \ill{} les ateliers \uncertain{qu.\ ens.}\ \uncertain{spécieux} resp.\ polyédraux plus généraux, il faut introduire sur $\mathcal{M}$ la notion de \struck{\ill{}} partie (fermée) « admissible » par des conditions convenables de « local finitude ».
\note{page de lecture très difficile, écrite vite ; beaucoup de mots n'y sont pas lus. L'articulation générale se laisse suivre : les « lieux » et les points, les espaces « infiniment divisibles », la structure « commune » aux contextes « spécieux » et « polyédral », les quatre axiomes Prat quens de la page 59, puis, pour des ateliers quasi-ensemblistes plus généraux, une notion de partie fermée « admissible » définie par des conditions de « finitude locale ». « commune » et « plus » sont soulignés sur la page}
\marginal{p.\ ex.\ \ill{} « Préatelier » \ill{} \ill{} \ill{} « Stratologie » (\ill{} « stratologique ») \ill{} \ill{} \ill{} \ill{} \ill{} \ill{} \ill{} topologie !}
\note{longue note oblique dans la marge inférieure gauche, en face de la seconde moitié de la page ; on y lit les mots « Préatelier », « Stratologie », « stratologique » et « topologie ! », le reste n'est pas lu}

\end{document}
